\documentclass[10pt, a4paper]{article}
\usepackage[margin=2cm]{geometry}
\usepackage[utf8]{inputenc}
\usepackage[portuguese]{babel}
\usepackage{amsmath}
\usepackage{amssymb}
\usepackage{mathtools}
\DeclareMathOperator{\Var}{Var}
\newcommand{\E}{\mathbb{E}}
\newcommand{\prob}{\mathbb{P}}

\title{\textbf{Lista 1 ME323}}
\author{Julio da Silva Telles RA:281275}

\begin{document}

\maketitle{}
\begin{enumerate}
    \item 
    \begin{enumerate}
        \item Determinístico
        \item Aleatório
        \item Determinístico
        \item Aleatório
    \end{enumerate}
    
    \item $\Omega = \{\text{livre}, \text{ocupada}, \text{em manutenção}\}$
    
    \item Considere dois conjuntos M e F sendo $M = \{\text{O (Onibus)}\text{, BI (Bicicleta)}\text{, CA (Carro)}\}$ os modos de se deslocar a faculdade e 
    $F = \{\text{C (Curta)} \text{, L (Longa)}\}$ os possiveis estados que a fila pode assumir temos: $$\Omega = M \times D$$ tal que 
    $\Omega = \{(\text{O},\text{C}),(\text{O},\text{L}),(\text{BI},\text{C}),(\text{BI},\text{L}),(\text{CA},\text{C}),(\text{CA},\text{L})\}$
    
    \item $A = \{\text{O}\} \text{, } B = \{\text{C}\} \text{, } A \cup B = \{\text{O ou C}\} \text{, } A \cap B = \{\text{O e C}\} \text{, } 
    A^{c} = \{\text{BI e CA}\} \text{, } A - B = A \cap B^{c} = \{\text{O e L}\}$
    
    \item Tome $\mathbb{L}$, $\mathbb{S}$ e $\mathbb{P}$ conjuntos que definem respectivamente a quantidade de linguagens de programação, sensores e protocolos
    disponiveis respectivamente. Observa-se que \#$\mathbb{L}$ = 3, \#$\mathbb{S}$ = 4, \#$\mathbb{P}$ = 2 temos então que pelo PFC (princípio fundamental da contagem): 
    
    $$\#\Omega = (\#\mathbb{L} * \#\mathbb{S} * \#\mathbb{P}) \Longrightarrow  \#\Omega = 24$$.          

    \item O numero de diferentes ordens que pode se ter de apresentações entre os alunos e $P_6 = 6!$ .
    
    \item  O numero de diferentes maneiras que o professor pode escolher os representantes da turma é:
    $$C_{8, 3} = P(X=k) = \binom{n}{k} p^k (1-p)^{n-k} $$ com $C_{8,3} = 56$.

    \item Expandindo $(x + y)^{4}$, temos: 
    $$(x + y)^{4} = x^{4} + 4x^{3}y + 6x^{2}y^{2} + 4y^{3}x + y^{4}$$ 
    tal que o coeficiente de $x^{2}y^{2} \text{, } C_1 = 6$.

    \item Tome $\mathbb{S}$, $\mathbb{C}$ e $\mathbb{M}$ conjuntos que definem respectivamente a quantidade de linguagens de programação, sensores e protocolos
    disponiveis respectivamente. Observa-se que \#$\mathbb{L}$ = 3, \#$\mathbb{S}$ = 2, \#$\mathbb{P}$ = 1 temos então que pelo PFC (princípio fundamental da contagem): 
    
    $$\#\Omega = (\#\mathbb{L} * \#\mathbb{S} * \#\mathbb{P}) \Longrightarrow  \#\Omega = 6$$. 

    \item Para cada um dos 4 minicursos, a estudante possui 2 opções (assistir ou não assistir). Pelo princípio multiplicativo, o total de subconjuntos possíveis é:
    $$ 2^4 = 16 $$

    \noindent Verificando através da soma das combinações:
    $$
    \begin{aligned}
    \text{Total} &= \binom{4}{0} + \binom{4}{1} + \binom{4}{2} + \binom{4}{3} + \binom{4}{4} \\
    &= 1 + 4 + 6 + 4 + 1 \\
    &= 16
    \end{aligned}
    $$

    \noindent Portanto, a estudante pode escolher entre \textbf{16} subconjuntos de minicursos. 

\end{enumerate} 

\end{document}  